\begin{description}
\item[(a)] Consider small surface element ($\Delta S$) of the disk. The
  gravitational field ($\Delta E$) generated by this surface element at at a
  point $O$ located at a distance $a$ from $\Delta S$ (see the figure below): % sic: "at at"
  \[
    \Delta E = G\frac{\sigma\,\Delta S}{a^2}
  \]
  where $\sigma$ is a surface density of the mass.

% FIGURE NEEDED: sketch of an irregular grey disk of surface density sigma with
% a small element dS, its surface normal, the angle alpha between the normal and
% the line to an external point O, and the solid angle dOmega subtended at O
% (2022 theory solutions, page 8 of 21).

  where $\alpha$ is the angle between surface normal and the vector $\vec r$
  and $\Delta\Omega$ is infinitesimal solid angle corresponding to the area
  $\Delta S$.

  For a ring, the surface density will be given by
  $\sigma = M/\pi(R^2 - r^2)$.

  As the horizontal component will get cancelled due to symmetry
  considerations, we only consider the normal component:
  \[
    \Delta E_\perp = \Delta E\cos(\alpha)
      = \frac{G\sigma}{a^2}\,\Delta S\cos(\alpha)
      = \frac{G\sigma}{a^2}\,\Delta S_\perp
      = G\sigma\,\Delta\Omega
  \]
  \hfill(5.0 points)

  By Summation, we can obtain total normal field of the surface:
  \[
    E_\perp = G\sigma\Omega
  \]
  \hfill(1.0 points)

  here $\Omega$ is the total solid angle subtended by the disk. In our case,
  when we observe disk from the point of view of the test particle:
  \begin{align*}
    \Omega &= 2\pi\left(1 - \cos\theta_{out}\right)
              - 2\pi\left(1 - \cos\theta_{in}\right) \\
           &= 2\pi\left(1 - \frac{x}{\sqrt{x^2 + R^2}}\right)
              - 2\pi\left(1 - \frac{x}{\sqrt{x^2 + r^2}}\right) \\
    \Omega &= 2\pi x\left(\frac{1}{\sqrt{x^2 + r^2}}
              - \frac{1}{\sqrt{x^2 + R^2}}\right)
  \end{align*}
  \hfill(3.0 points)
  \begin{align*}
    \therefore F_\perp &= G\sigma\Omega m \\
    F_\perp &= \frac{2GMmx}{R^2 - r^2}
       \left(\frac{1}{\sqrt{x^2 + r^2}} - \frac{1}{\sqrt{x^2 + R^2}}\right)
  \end{align*}
  \hfill(1.0 points)

\item[(b)] In small displacement approximation:
  \[
    F_\perp \approx \frac{2GMmx}{R^2 - r^2}
      \left(\frac{1}{r} - \frac{1}{R}\right)
      = \left(\frac{2GMm}{Rr(R + r)}\right)x = kx
  \]
  \hfill(3.0 points)

  when the distance between the disk and the point mass is $x$, the distance
  of the center of mass of the system from the center of the disk will be:
  \[
    x_1 = \frac{mx}{m + M}
  \]
  \hfill(2.5 points)

  If we choose the center of mass of the system as an origin of our coordinate
  system, then equation of motion of the disk will have a form:
  \begin{align*}
    Ma_1 &= -kx \\
    Ma_1 + k\,\frac{m + M}{m}\,x_1 &= 0
  \end{align*}
  \hfill(2.5 points)

  For this simple harmonic oscillator, the period of small oscillations will
  be:
  \[
    T = 2\pi\sqrt{\frac{Rr(R + r)}{2G(M + m)}}
  \]
  \hfill(2.0 points)
\end{description}
